\documentclass[a4paper]{article}
\usepackage{graphicx} % Required for inserting images
\usepackage{wrapfig}
\usepackage{enumitem}
\usepackage[center]{caption}
\usepackage[english, russian]{babel}
\usepackage{amsmath}
\usepackage{geometry}
\geometry{verbose, a4paper, tmargin=2cm, bmargin=2cm, lmargin=2.0cm, rmargin=2.0cm}

\begin{document}

\begin{titlepage}
	\centering
	{\scshape\Large МОСКОВСКИЙ ФИЗИКО-ТЕХНИЧЕСКИЙ ИНСТИТУТ \\
	(НАЦИОНАЛЬНЫЙ ИССЛЕДОВАТЕЛЬСКИЙ УНИВЕРСИТЕТ)\\ % название ВУЗа большим шрифтом
    Физтех-школа аэрокосмических технологий}
	
	\vspace{4cm} % отступ 4 см по вертикали
	{\LARGE Отчет о выполнении лабораторной работы 1.2.1}
	
	\vspace{1cm} % ещё 1 см
	{\huge\bf Определение скорости полета пули при помощи баллистического маятника}
	
	\vspace{1cm} % ещё 1 см
	\vfill % заполнение по вертикали пустотой (LaTeX сам решит, сколько здесь отступить)
	
    \begin{flushright} % Этот текст будет с правого края страницы
	   {\LARGE Ефремова Татьяна, Б03-503}
    \end{flushright}
    \vfill

	\today % Дата сборки документа
\end{titlepage}

\section{Аннотация}
Цели работы: определить скорость полета пули, применяя законы сохранения и используя баллистические маятники.

\section{Теоретические сведения}
Поскольку скорость вылета пули из духового ружья обычно порядка $100 - 1000$ м/с, а размер лабораторной установки -- порядка нескольких метров, время полета пули составляет величину порядка $10^{-3}-10^{-2}$с. Таким образом, определить скорость пули легче не напрямую, а измерив импульс, передаваемый некоторому телу при неупругом соударении. \\ \\
В данной работе для этого будет использоваться баллистический маятник. Так как внешние силы отсутсвуют, а удар кратковременный (время соударения много меньше периода колебания маятника), можно считать, что импульс системы сохраняется. В дальнейшем будем считать затухание малым, если за десять колебаний амплитуда уменшается меньше, чем наполовину. В таком случае свзязь между максимальным отклонением маятника и начальной скорости можно описывать законом сохранения механической энергии. По начальному максимальному отклонению определяются импульс и скорость пули. \\ \\
При проведении эксперимента необходимо, чтобы колебания маятника происходили в одной плоскости, а воздушная струя, выходящая вслед за пулей из ружья, не оказывала влияния на его движение. Это достигается расположением ружья строго перпендикулярно плоскости мишени маятника на достаточном расстоянии от нее.

\subsection{Баллистический маятник, совершающий поступательное движение}

Используемый в этой части работы баллистический маятника представлет собой тяжелый циллиндр, подвешенный на четырех нитях одинаковой длины. Он изображен на рис. 1 вместе с измерительной системой. \\ 
\begin{figure}[!h]
   \centering    
    \includegraphics[width=0.7\linewidth]{C:/Users/User/Documents/mipt_takonef/лабы/1.2.1/1.2.1_1.png}
    \caption{Схема установки для измерения скорости полета пули} 
\end{figure} \\
Горизонтальные компоненты внешних сил (тяжести и натяжения нитей) либо малы (т. к. малы отклонения), либо вовсе отсутствуют. Тем более мал по сравнению с импульсом пули их импульс во время соударения. Тогда закон сохранения импульса при соударении пули с цилиндром имеет вид
\begin{equation}
mu = (M+m)V,
\end{equation}
где $m$ -- масса пули, $M$ -- масса циллиндра, $V$ -- скорость циллиндра и пули после неупругого столкновения. \\\\
Учитывая, что масса маятника значительно больше массы пули,
\begin{equation}
u = \frac{M}{m}V.
\end{equation}
Принебрегая потерями, по закону сохранения механической энергии
\begin{equation}
V^2 = 2gh,
\end{equation}
где $h$ -- высота подъема маятника над его изначальным положением, $g$ -- ускорение свободного падения. \\\\
Высота подъема маятника выражается через угол $\varphi$ отклонения нитей подвеса от вертикали:
\begin{equation}
h = L(1 - cos\varphi) = 2Lsin^2\frac{\varphi}{2}, \quad \varphi \approx \frac{\Delta x}{L},
\end{equation}
где $L$ -- длина нити подвеса, $\Delta x$ -- отклонение маятника по горизонтали. \\\\
Тогда из уравнений (1) - (4):
\begin{equation}
u = \frac{M}{m} \sqrt{\frac{g}{L}} \Delta x.
\end{equation}
Измерение $\Delta x$ производится с помощью оптической системы (рис. 1) по увеличенной горизонтальной шкале на экране. 

\subsection{Крутильный баллистический маятник}
Схема установки изображена на рис. 2. Пуля массой $m$ попадает в мишень, укрепленную на стержне $a$, который вместе с грузами $M$ и проволокой П образует крутильный маятник. \\
\begin{figure}[!h]
   \centering    
    \includegraphics[width=0.67\linewidth]{C:/Users/User/Documents/mipt_takonef/лабы/1.2.1/1.2.1_2.png}
    \caption{Схема установки для измерения скорости полета пули} 
\end{figure} \\
Для определения скорости $u$ полета пули непосредственно перед ударом воспользуемся законом сохранения момента импульса в виде 
\begin{equation}
mur = I\Omega,
\end{equation}
где $r$ -- расстояние от линии полета пули до оси вращения маятника (до проволоки П), $I$ -- момент инерции маятника, $\Omega$ -- его угловая скорость непосредственно после удара.\\ \\
Пренебрегая потерями, по закону сохранения механической энергии
\begin{equation}
k\frac{\varphi^2}{2}  = I\frac{\Omega^2}{2},
\end{equation}
где $k$ -- модуль упругости кручения проволоки, $\varphi$ -- максимальный угошл поворота маятника.\\\\
Из (6) и (7):
\begin{equation}
u = \varphi \frac{\sqrt{kl}}{mr}.
\end{equation}
Т. к. угол закручивания $\varphi$ мал, его можно найти как
\begin{equation}
\varphi \approx \frac{x}{d},
\end{equation} 
где $x$ -- смещение, $d$ -- расстояние от шкалы до оси вращения. \\\\
Период колебания крутильного маятника без грузов
\begin{equation}
T_1 = 2\pi \sqrt{\frac{I}{k}}.
\end{equation}
Период колебания крутильного маятника с грузами $M$.
\begin{equation}
T_2 = 2\pi \sqrt{\frac{I-2MR^2}{k}},
\end{equation}
Из (10) и (11):
\begin{equation}
\sqrt{kI}  = \frac{4\pi MR^2 T_1}{T_1^2 - T_2^2}
\end{equation}
где $R$ -- расстояние от центров грузов до проволоки.

\section{Оборудование}
\subsection{Используемое оборудование}
Духовое ружье на штативе, осветлитель, оптическая система для измерения отклонений маятника, измерительная линейка, пули и весы для их взвешивания, а также баллистические маятники. 
\subsection{Инструментальные погрешности}
\vspace{0.15 cm}

\setdescription{leftmargin=0pt}
\begin{description}
    \item[линейка:] $\Delta_\text{лин}=\pm0,5$ мм.
    \item[весы:] $\Delta_\text{в}=\pm 5$ мг.
    \item[оптическая система:] $\Delta_\text{ОС}=\pm0,5$ мм.
\end{description}

\section{Результаты измерений и обработка данных}
Масса циллиндра поступательного маятника $M = 2900 \pm 5$ г. Длина нити подвеса $L = 2170 \pm 5$ мм.

\begin{table}[!h]
\centering
\caption{Амплитуды колебаний на шкале, скорости пуль}
\begin{tabular}{|c|c|c|c|}        \hline
    №Пули & m, мг & $\Delta x$, мм & $u$, м/с \\  \hline
    1 & $509$ & $11.6$ & $140$\\ \hline
    2 & $508$ & $11.5$ & $139$\\ \hline
    3 & $505$ & $11.4$ & $138$\\ \hline
    4 & $500$ & $11.4$ & $138$\\ \hline
    5 & $510$ & $11.5$ & $139$\\ \hline
\end{tabular}
\end{table}

\newpage
Рассчитаем начальную скорость пули по формуле:

\[u = \frac{M}{m} \sqrt{\frac{g}{L}} \Delta x\]

Рассчитаем погрешность скорости:

\[\varepsilon_{u_{\text{снст}}} = \varepsilon_M + \varepsilon_m + \frac{1}{2} \varepsilon_g + \frac{1}{2} \varepsilon_L + \varepsilon_{\Delta x} \approx 7.5\%\]

\[\sigma_{u_{\text{снст}}} = u \varepsilon_{u_{\text{снст}}} \approx 10.4 \text{м/с}\]

\[\sigma_{u_{\text{случ}}} \approx 5.6 \text{м/с}\]

\[\sigma_u = \sqrt{\sigma_{u_{\text{снст}}}^2 + \sigma_{u_{\text{случ}}}^2} \approx 11.8 \text{м/с} \]

\[\varepsilon_u = \frac{\sigma_u}{u} \approx 8.5\%\]

Тогда:

\[u = (139 \pm 12) \text{м/с}\]
\\
Измерим параметры крутильного баллистического маятника (рис. 2):
\begin{table}[!h]
\centering
\caption{Характеристики крутильного маятника}
\begin{tabular}{|c|c|c|c|c|}        \hline
    $d$, мм & $R$, мм & $r$, мм & $M_{g1}$, г& $M_{g2}$, г\\ \hline
    1387 & 336 & 213 & 725.6 & 724.5 \\ \hline
\end{tabular}
\end{table} \\
Результаты измерения времени одного полного колебания маятника (по результатам 10 колебаний):
\begin{table}[!h]
\centering
\caption{ Периоды колебаний крутильного маятника}
\begin{tabular}{|c|c|}        \hline
    Измерение & Период, с \\  \hline
    Без грузов, $T_1$ & $5.2\pm 0.1$  \\ \hline
    С грузами, $T_2$ & $6.9\pm 0.1$ \\ \hline
\end{tabular}
\end{table} \\
Отсюда:
\[\sqrt{kI} = \frac{4\pi MR^2 T_1}{T_1^2 - T_2^2} \approx (0.699 \pm 0.006)^{\frac{\text{кг} \cdot \text{м}^2}{\text{с}}}\]

Результаты измерения горизонтального смещения крутильного маятника на шкале, скорость пули для данной амплитуды.
\begin{table}[!h]
\centering
\caption{Амплитуды колебаний на шкале, скорости пуль}
\begin{tabular}{|c|c|c|c|}        \hline
    №Пули &  m, мг & $\Delta x$, мм & $u$, м/c \\  \hline
    6 & $510 \pm 5$ & $42 \pm 2$ & 96$\pm 9$\\ \hline
    7 & $512 \pm 5$ & $39 \pm 2$ & 99$\pm 9$\\ \hline
    8 & $509 \pm 5$ & $49\pm 2$ & 100$\pm 9$\\ \hline
    9 & $507 \pm 5$ & $45\pm 2$ & 102$\pm 9$\\ \hline
    10 & $513 \pm 5$ & $46\pm 2$ & 102$\pm 9$\\ \hline
\end{tabular}
\end{table} \\
Рассчитаем начальную скорость пули по формуле:

\[u = \varphi \frac{\sqrt{kI}}{mr} \approx x \frac{\sqrt{kI}}{2dmr}\]

\[\varepsilon_u = \varepsilon_{\sqrt{kI}} + \varepsilon_x + \varepsilon_m + \varepsilon_r \approx 5.8\%\]

Рассчитаем погрешность:

\[\sigma_{u_{\text{случ}}} \approx 10^{\frac{\text{м}}{\text{с}}}\]

\[\sigma_{u_{\text{сист}}} = u \varepsilon_u \approx 5^{\frac{\text{м}}{\text{с}}}\]

\[\sigma_u = \sqrt{\sigma_{u_{\text{сист}}}^2 + \sigma_{u_{\text{случ}}}^2} \approx 10.9^{\frac{\text{м}}{\text{с}}}\]

По итогу:

\[u = (94 \pm 11)^{\frac{\text{м}}{\text{с}}} (11.7\%)\]


Среднее значение и разброс:

Для поступательного маятника средняя скорость пули $u_{mid}=126\pm 1$ м/с

Для крутильного маятника средняя скорость пули $u_{mid}=100\pm 1$ м/с

\section{Выводы}


\end{document}
